\chapter{Discusión}
\label{cap:discusion}

\section{Qué demuestra el sistema}

El aporte principal de ManyHands no es generar más líneas de código ni reemplazar
la capacidad de los modelos utilizados como ejecutores. Su aporte es una cadena
de custodia para el trabajo generado:

\begin{equation}
\text{Objetivo} \rightarrow \text{Plan} \rightarrow \text{Contratos}
\rightarrow \text{Intentos} \rightarrow \text{Evidencia}
\rightarrow \text{Integración} \rightarrow \text{Entrega}.
\label{eq:cadena-custodia}
\end{equation}

El run Viaje en Familia muestra que esta cadena puede operar de extremo a
extremo. La planificación produjo una jerarquía coherente; el scheduler encontró
fronteras concurrentes; los compuestos integraron a sus hijos; la recuperación
reabrió nodos fallidos sin descartar artefactos adoptados; y la entrega quedó
vinculada a un commit y un árbol Git reproducibles.

El resultado no implica que toda ejecución autónoma de software vaya a
completar cualquier objetivo. Demuestra una proposición más acotada y útil: es
posible construir un plano de control que mantenga trazabilidad y permita
recuperar fallas durante una colaboración multi-agente real.

\section{Modelos para semántica, software para autoridad}

Una decisión transversal del diseño fue separar dos responsabilidades:

\begin{itemize}
  \item el modelo interpreta el objetivo, propone límites semánticos y escribe
  o repara código;
  \item el plano determinista valida esquemas, calcula disponibilidad, limita
  recursos, inspecciona Git y decide si existe evidencia suficiente.
\end{itemize}

La separación evita pedir al mismo componente probabilístico que produzca el
resultado y, al mismo tiempo, sea la autoridad final sobre su corrección. El
modelo conserva libertad dentro de un intento; los contratos y verificadores
limitan las consecuencias que ese intento puede tener sobre el run.

\section{Por qué la granularidad es categórica}

La experiencia con el puntaje continuo mostró que una fórmula puede parecer
precisa sin ser explicable. Sumar cantidad de archivos, tokens, interfaces y
pruebas produce un número, pero no una razón física para dividir. La política
categórica cambia la pregunta: en lugar de ``¿qué puntaje obtuvo esta tarea?'',
pregunta ``¿qué beneficio concreto compra la separación?''.

Las tres respuestas posibles —capacidad, paralelismo o aislamiento de
validación— pueden explicarse a un operador y observarse después de la
ejecución. Además, el algoritmo de colapso reduce el costo de una propuesta
demasiado fragmentada sin perder los criterios que debían verificarse.

\section{Costo de control y recuperación}

Los worktrees, journals, validaciones y recibos agregan tiempo respecto de una
invocación única. Este costo es deliberado. En un parche trivial puede ser
innecesario; en un cambio con varias fronteras, permite atribuir resultados,
aislar fallas y continuar después de una interrupción.

El experimento también muestra que la recuperación tiene costo propio. Un nodo
puede requerir varios intentos antes de converger, y una evidencia mal
especificada puede bloquear código correcto. La ventaja no es eliminar el
fallo, sino volverlo local, durable y diagnosticable. La principal oportunidad
de mejora consiste en reducir intentos repetidos mediante mejores diagnósticos
y checkpoints de candidatos parciales.

\section{Alcance de la evidencia}

La evidencia técnica es fuerte en aquellas propiedades vinculadas a Git y al
journal: topología, orden de eventos, procesos, candidatos, matrices, entrega y
reproducción. Las 32 pruebas del target cubren la lógica de las áreas y su
composición. La revisión visual confirma la presentación del dashboard y sus
estados básicos.

La evaluación no constituye un estudio de usabilidad con usuarios ni una
comparación estadística contra otros orquestadores. Tampoco demuestra
escalabilidad a repositorios de millones de líneas o a todos los lenguajes. Estas
distinciones protegen la interpretación del caso: el experimento aporta
evidencia de viabilidad arquitectónica, no una ley universal sobre rendimiento
de modelos.

\section{Amenazas a la validez}

\begin{itemize}
  \item \textbf{Variabilidad del modelo.} Una nueva ejecución puede proponer
  código o particiones diferentes. El freeze de configuración y el journal
  reducen ambigüedad, pero no eliminan la estocasticidad.
  \item \textbf{Caso único.} Viaje en Familia fue elegido por su riqueza
  estructural. Otros dominios pueden producir grafos y costos distintos.
  \item \textbf{Pruebas como aproximación.} Una suite verde no mide por sí sola
  mantenibilidad o experiencia de uso. Por eso se separan las conclusiones de
  ingeniería de las futuras evaluaciones cualitativas.
  \item \textbf{Entorno local.} La supervisión y el sandbox fueron ejercitados
  en Windows. La portabilidad conceptual no sustituye una calificación física
  equivalente en otros sistemas operativos.
\end{itemize}

Estas amenazas no invalidan el recorrido observado, pero delimitan qué puede
inferirse a partir de él.
